\documentclass[12pt,a4paper]{report}
\usepackage[utf8]{inputenc}
\usepackage[english]{babel}
\usepackage{amsmath}
\usepackage{amsfonts}
\usepackage{amssymb}
\usepackage{graphicx}
\usepackage{hyperref}
\usepackage{booktabs}
\usepackage{float}
\usepackage{geometry}
\usepackage{setspace}
\usepackage{caption}
\usepackage{subcaption}
\usepackage{titlesec}
\usepackage[T1]{fontenc}
\usepackage[scaled=0.92]{helvet} 
\renewcommand{\familydefault}{\sfdefault}

% Aligning rigorously with the ENSSEA standard
\geometry{margin=2.5cm, left=3cm} 
\setstretch{1.5}

\titleformat{\chapter}[display]
  {\normalfont\huge\bfseries\centering}{\chaptertitlename\ \thechapter}{20pt}{\Huge}
\titleformat{\section}
  {\normalfont\Large\bfseries}{\thesection \ - }{1em}{}
\titleformat{\subsection}
  {\normalfont\large\bfseries}{\thesubsection \ - }{1em}{}
\titleformat{\subsubsection}
  {\normalfont\normalsize\bfseries\underline}{\thesubsubsection \ - }{1em}{}

\begin{document}

\setcounter{chapter}{1} % Set chapter to 2
\chapter{RELATIONSHIPS BETWEEN VARIABLES, LITERATURE REVIEW AND INTERNSHIP FRAMEWORK}

\section*{Introduction of Chapter II}
\addcontentsline{toc}{section}{Introduction of Chapter II}
The transition from generalized, archaic marketing paradigms towards highly targeted, algorithmic enterprise intelligence inherently demands a rigorous theoretical foundation. Prior to deploying empirical machine learning infrastructures—as will be demonstrated profoundly within the structural data of Chapter III—it is fundamentally imperative to establish the strictly defined epistemological boundaries linking the core components of the study. Algorithmic Customer Relationship Management (CRM) does not function as an isolated set of disjointed calculations; it operates as a continuous, interlocking matrix where the outputs of one psychological boundary directly fuel the probabilities of another.

This chapter is entirely devoted to establishing the conceptual framework dictating the profound interaction existing between the study’s primary independent and dependent variables. Specifically, Section 2.1 will meticulously dissect the theoretical literature governing the causal relationship looping Customer Segmentation (RFM \& K-Means) directly into the probabilistic generation of Customer Lifetime Value (CLV). Furthermore, it bridges the micro-economic behavior of individual retail actors into the macro-economic reality of aggregate Sales Forecasting (SARIMA, Prophet, XGBoost). Concluding this section, a synthesized conceptual model is constructed, formulating the exact five research hypotheses (H1–H5) that the subsequent empirical pipeline will mathematically either validate or definitively reject in the context of the Algerian retail sector.

% -----------------------------------------------------------------------------
\section{Relationships between the Variables of the Study}

\subsection{Relationship between Customer Segmentation and Customer Lifetime Value (CLV)}
Customer Lifetime Value (CLV)—defined traditionally within marketing literature as the discounted sum of all future net cash flows generated by an individual consumer—represents the absolute apex of CRM mechanics\footnote{Blattberg, R. C., Getz, G., \& Thomas, J. S. : \textit{Customer Equity: Building and Managing Relationships As Valuable Assets}, Harvard Business School Press, Boston, 2001, p. 88.}. However, mathematical probability models calculating forward-looking yield are notoriously unstable and entirely inaccurate when forcefully applied across a heterogeneous, globally generalized audience. 

The structural relationship between Segmentation and CLV is fundamentally chronological and causal. Statistical probability fields, such as the Beta-Geometric Negative Binomial Distribution (BG/NBD) governing churn, demand behavioral homogeneity internally to prevent outlier vectors from completely rupturing the Gamma distribution parameters\footnote{Fader, P. S., \& Hardie, B. G. S. : \textit{Probability Models for Customer-Base Analysis}, Journal of Interactive Marketing, 23(1), 2009, p. 61.}. Consequently, preliminary heuristic segmentation—specifically the extraction of Recency, Frequency, and Monetary (RFM) components alongside unsupervised K-Means Euclidean partitions—acts as the mandatory precursor variable.

\begin{itemize}
    \item \textbf{Isolation of the Actionable Core}: Continuous tracking studies assert that practically $60\%$ to $70\%$ of a physical retail database constitutes "One-Time" or strictly opportunistic actors. Attempting to formulate probabilistic lifespan tracking on single-purchase actors creates catastrophic arithmetic denominators. Rigorous RFM Segmentation explicitly isolates "Champions" and "Loyal Customers," fundamentally validating the sub-cohorts capable of executing the fractional sub-variables required for Gamma-Gamma average margin tracking.
    \item \textbf{Derivation of the Churn Horizon}: The mathematical definition of a "lost" customer diverges wildly depending on their cluster. A localized "Champion" suddenly missing for 60 days inherently spikes the probability of defection ($P(\text{Alive})$ dropping significantly), whereas a "Dormant" buyer vanishing for 60 days strictly aligns with their historical pattern constraint. Without spatial clustering variables executed concurrently, algorithmic CLV equations inherently remain blind to varying psychological tolerance limits.
\end{itemize}
In summation, Customer Segmentation is the theoretical independent variable that mathematically stabilizes the dependent execution of Customer Lifetime Value parameters.

\subsection{Relationship between CLV and Sales Forecasting: Analytical Interdependencies}
While CLV models evaluate the micro-level lifespan and fractional value of individual human behavioral actors, Sales Forecasting models (such as SARIMA and XGBoost) calculate the macro-level volumetric inventory requirements functioning universally across physical enterprise logistics\footnote{Makridakis, S., Wheelwright, S. C., \& Hyndman, R. J. : \textit{Forecasting: Methods and Applications}, John Wiley \& Sons, New York, 1998, p. 24.}. Within archaic corporate operations, these two dimensions are notoriously structured entirely separate from each other, housed within distinct departmental silos (Marketing versus Logistical Supply Chain).

Modern theoretical frameworks posit that the relationship bridging CLV intimately against aggregate Forecasting is fundamentally interdependent. Macro-level prediction tools dictate \textit{what} logistical volume is expected and \textit{when} the violent seasonal spikes will physically occur. Conversely, the CLV metric identifies precisely \textit{who} is structurally generating that base volume capacity. 

The academic correlation tracking between the variables operates effectively as an acquisition damper mechanism:
\begin{enumerate}
    \item \textbf{Baseline Stabilization via High CLV}: A retail network possessing localized clusters harboring high continuous Median Lifetime Values fundamentally exhibits superior stationary demand curves. High loyalist volume mathematically suppresses extreme time-series variance ($ \epsilon_t $), generating deeply stabilized autoregressive mapping frameworks which allow algorithms like SARIMA to effortlessly anticipate continuous structural baseline cash flows.
    \item \textbf{Forecasting as the Operational Damper}: While probabilistic CLV identifies customers fundamentally maintaining high structural Net Present Value (NPV), this capacity cannot physically manifest functionally if logistical operations incur an "Out-of-Stock" crisis. Precise aggregated temporal forecasting mathematically ensures targeted "Champions" arriving at physical establishments during predicted seasonal weekend spikes natively find inventory completely staged and ready. A failure within the Forecasting variable generates structural supply chain collapse, directly destroying continuous CLV generation execution natively.
\end{enumerate}
Thus, advanced logistics literature verifies that predicting long-term customer relationships requires simultaneous inventory time-series assurance to complete the systemic transactional cycle.

\subsection{Integrated Conceptual Model and Formulation of Research Hypotheses (H1-H5)}
Drawing profoundly upon the intricate mathematical interactions established structurally between heuristic behavior limits, unsupervised spatial boundaries, probabilistic tracking, and macro-aggregate time prediction, an integrated conceptual model natively targeting modern Algerian retail environments can logically be synthesized.

This conceptual framework fundamentally shifts logistical operations entirely away from instinctual estimation, anchoring the core of the PUMA commercial enterprise structure upon strict algorithmic probability execution. To critically validate this specific conceptual integration empirically via Python 3 mechanics in the subsequent Chapter III, five rigorous hypotheses are formally established:

\begin{itemize}
    \item \textbf{Hypothesis 1 (H1)}: Executing strict categorical heuristic RFM rules coupled physically alongside unsupervised continuous log-normalized K-Means geometry isolates specific, statistically predictable behavioral cohorts significantly better than generalized mass audience tracking constraints.
    \item \textbf{Hypothesis 2 (H2)}: Consumer segments algebraically defined via heuristic parameters as "Champions" and "Loyal Customers" inherently display statistically significant higher Median Customer Lifetime Value (CLV) yield projections than cohorts localized within the "Promising" bounds.
    \item \textbf{Hypothesis 3 (H3)}: Operating exclusively within physically brutal, highly turbulent Algerian retail environments defined strictly by violent weekend oscillations, Autoregressive matrices mathematically remembering past sequential errors (SARIMA) comprehensively outperform linear, additive trend frameworks (Prophet) in decreasing continuous absolute volumetric errors natively.
    \item \textbf{Hypothesis 4 (H4)}: Implementing Beta-Geometric decay methodologies specifically accurately flags transitioning "At Risk" actors, isolating precise $P(\text{Alive})$ failure thresholds mathematically prior to totally invisible defection materialization.
    \item \textbf{Hypothesis 5 (H5)}: Constructing an algorithmic, unified interactive Executive Dashboard mathematically overlapping aggregate Sales Forecasting limits intimately against individual CLV parameters structurally generates a fundamentally robust enterprise framework capable of drastically lowering operational Logistical overheads.
\end{itemize}

These explicit hypotheses structurally bridge the theoretical literature directly into the physical computational pipelines constructed within the PUMA historical ecosystem, formally defining the rigorous boundaries verified conclusively within the empirical application stages.

\end{document}
